\chapterimage{figs/chapterhead.png}
\chapter{Repository Architecture}
\label{chap:repository-architecture}

This chapter documents the current repository layout as it exists in the active
branch and build surface. It is based on the root \texttt{README.md}, the root
\texttt{CMakeLists.txt}, \texttt{CMakePresets.json}, the canonical
architecture/status documents, and the filesystem tree observed on
2026-07-23. It does not treat historical paths or planned workspaces as if they
were current build boundaries.

\noindent\textbf{Objective.} Explain how the repository is organized, which
directories are source architecture, and how the current CMake hierarchy maps to
real modules, targets, executables, tests, and support areas.

\noindent\textbf{Audience.} Developers who need a faithful map of the codebase
before changing modules, build wiring, documentation, or packaging.

\noindent\textbf{Scope.} Cover the top-level repository tree, the source tree,
the current CMake target families, the non-build support areas, and the current
reading path for technical work.

\noindent\textbf{Status.} Confirmed by repository inspection on 2026-07-23.
The chapter describes the current architecture snapshot, not a future idealized
layout.

The chapter follows this sequence:
Figure~\ref{fig:repo-tree},
Figure~\ref{fig:repo-target-graph},
Figure~\ref{fig:repo-layer-map},
Figure~\ref{fig:repo-boundary-map}, and
Figure~\ref{fig:repo-build-flow}.

\section{Repository snapshot}
\label{sec:repo-snapshot}

The repository is intentionally organized around a small number of top-level
roots:
\begin{itemize}
\item \texttt{CMakeLists.txt} and \texttt{CMakePresets.json} define the build
surface;
\item \texttt{source/} contains the compiled code and test code;
\item \texttt{models/} contains example and didactic \texttt{.gen} inputs;
\item \texttt{docs/} contains the manuals and governance material;
\item \texttt{debian/}, \texttt{packaging/}, and \texttt{scripts/} contain
packaging and automation helpers.
\end{itemize}

Generated build directories such as \texttt{build/} or
\texttt{cmake-build-debug/} exist in the checkout, but they are not source
architecture and must not be read as canonical repository structure.

\begin{figure}[htbp]
\centering
% FIGURE-SPEC-BEGIN
% ID: fig:repo-tree
% Source of truth:
%   - ../README.md
%   - ../docs/ai_assistants/ARCHITECTURE.md
%   - ../CMakeLists.txt
%   - filesystem tree under ../
% Purpose:
%   Show the current top-level repository layout and the major source/support roots.
% Required visual content:
%   - top-level roots
%   - source subtree families
%   - docs, models, packaging, scripts, debian
% Accessibility description:
%   A compact repository map that helps the reader orient before drilling into CMake or source trees.
% Validation:
%   Every directory label must exist in the current checkout and the diagram must not include generated build directories as source architecture.
% FIGURE-SPEC-END
\figureplaceholder{figs/generated/repository-tree}
\caption{Current repository tree and the major top-level roots.}
\label{fig:repo-tree}
\end{figure}

Figure~\ref{fig:repo-tree} should make the boundaries obvious at a glance: the
code lives under \texttt{source/}, the reference models live under
\texttt{models/}, and the documentation and packaging helpers are separate
subtrees rather than hidden inside the application code.

\section{Source tree and build graph}
\label{sec:repo-source-tree}

The current build graph is rooted in \texttt{CMakeLists.txt}. That file does not
build the repository by directory name alone; it turns on and off module groups
through explicit options and then adds only the relevant subdirectories.

The current root options are:
\begin{itemize}
\item \texttt{GENESYS\_BUILD\_KERNEL};
\item \texttt{GENESYS\_BUILD\_PARSER};
\item \texttt{GENESYS\_BUILD\_PLUGINS};
\item \texttt{GENESYS\_BUILD\_TERMINAL\_APPLICATION};
\item \texttt{GENESYS\_BUILD\_TERMINAL\_EXAMPLES};
\item \texttt{GENESYS\_BUILD\_WORKER\_APPLICATION};
\item \texttt{GENESYS\_BUILD\_GUI\_APPLICATION};
\item \texttt{GENESYS\_BUILD\_TOOLS};
\item \texttt{GENESYS\_BUILD\_TESTS};
\item \texttt{GENESYS\_ENABLE\_PYTHON\_INTEGRATION}.
\end{itemize}

The compatibility aliases currently retained in the build layer are:
\begin{itemize}
\item \texttt{GENESYS\_TERMINAL\_APPLICATION};
\item \texttt{GENESYS\_BUILD\_WEB\_APPLICATION};
\item \texttt{GENESYS\_TERMINAL\_EXAMPLE};
\item \texttt{GENESYS\_TERMINAL\_EXAMPLE\_CLASS}.
\end{itemize}
They are transitional names only. The planned GUI experiment slot is
\texttt{GENESYS\_BUILD\_GUI\_DOEXPERIMENTS}, and it is intentionally guarded
as not-yet-implemented, so it is not treated here as a current architecture
surface.

\begin{figure}[htbp]
\centering
% FIGURE-SPEC-BEGIN
% ID: fig:repo-target-graph
% Source of truth:
%   - ../CMakeLists.txt
%   - ../CMakePresets.json
%   - ../source/kernel/*/CMakeLists.txt
%   - ../source/plugins/*/CMakeLists.txt
%   - ../source/parser/CMakeLists.txt
%   - ../source/applications/*/CMakeLists.txt
%   - ../source/tests/CMakeLists.txt
%   - ../source/tools/CMakeLists.txt
% Purpose:
%   Show the actual target families and their link/build dependencies.
% Required visual content:
%   - static core libraries
%   - application executables
%   - test executables and aggregate targets
%   - compatibility aliases and custom targets where they matter
% Accessibility description:
%   A dependency graph that helps the reader distinguish libraries, executables, aliases, and aggregate targets.
% Validation:
%   Every node must correspond to a real current target or alias, and every edge must follow a real target_link_libraries or add_dependencies relation.
% FIGURE-SPEC-END
\figureplaceholder{figs/generated/repository-target-graph}
\caption{Current target graph derived from the live CMake files.}
\label{fig:repo-target-graph}
\end{figure}

Figure~\ref{fig:repo-target-graph} is the canonical view of the current build
surface. The core static chain is the following dependency ladder:
\begin{itemize}
\item \texttt{genesys\_kernel\_util};
\item \texttt{genesys\_kernel\_statistics};
\item \texttt{genesys\_kernel\_simulator\_support};
\item \texttt{genesys\_kernel\_simulator\_runtime}.
\end{itemize}

The plugin and parser side feed the same runtime chain. The relevant static
libraries are:
\begin{itemize}
\item \texttt{genesys\_plugins\_data};
\item \texttt{genesys\_plugins\_components\_minimal};
\item \texttt{genesys\_plugins\_components} for the broader component
library with an optional GLPK path;
\item \texttt{genesys\_parser}, which remains a separate parser library with
optional Bison/Flex regeneration.
\end{itemize}

The executable families hang off that shared core:
\begin{itemize}
\item \texttt{genesys\_shell};
\item \texttt{genesys\_modelspecific\_app} and the legacy
\texttt{genesys\_terminal\_application} copy target;
\item \texttt{genesys\_worker\_app} / \texttt{genesys\_worker};
\item \texttt{genesys\_gui\_application} / \texttt{genesys\_gui};
\item \texttt{genesys\_httpworker\_gui\_application};
\item \texttt{genesys\_dataanalyser\_gui\_application};
\item \texttt{genesys\_optimizer\_gui\_application};
\item \texttt{genesys\_ai\_assistant\_gui\_application};
\item \texttt{genesys\_tools\_application};
\item the test executables under \texttt{source/tests/}.
\end{itemize}

\section{Layered boundaries}
\label{sec:repo-layered-boundaries}

The repository is easier to understand when read in layers rather than as a
flat directory listing.

\begin{figure}[htbp]
\centering
% FIGURE-SPEC-BEGIN
% ID: fig:repo-layer-map
% Source of truth:
%   - ../docs/ai_assistants/ARCHITECTURE.md
%   - ../README.md
%   - ../CMakeLists.txt
%   - ../source/kernel/*
%   - ../source/parser/*
%   - ../source/plugins/*
%   - ../source/applications/*
%   - ../source/tools/*
% Purpose:
%   Show the durable architectural layers and how source families map onto them.
% Required visual content:
%   - kernel as the foundation
%   - parser and plugins as language/extensibility layers
%   - applications and tools on top
%   - tests as validation
%   - non-source roots as support material
% Accessibility description:
%   A layer diagram that explains which parts of the tree are structural and which are support or content.
% Validation:
%   The layer names must match current repository boundaries and the diagram must not imply forbidden cross-layer ownership.
% FIGURE-SPEC-END
\figureplaceholder{figs/generated/repository-layer-map}
\caption{Current repository layers and their responsibilities.}
\label{fig:repo-layer-map}
\end{figure}

Figure~\ref{fig:repo-layer-map} captures the most useful reading order for new
work:
\begin{enumerate}
\item kernel;
\item parser and plugins;
\item applications and tools;
\item tests;
\item support material such as docs, models, packaging, and scripts.
\end{enumerate}

The kernel contains the runtime foundation:
\begin{itemize}
\item \texttt{source/kernel/util/};
\item \texttt{source/kernel/statistics/};
\item \texttt{source/kernel/simulator/}.
\end{itemize}
The parser owns expression interpretation in \texttt{source/parser/}. The
plugins layer owns data and component definitions in
\texttt{source/plugins/data/} and \texttt{source/plugins/components/}.

The applications layer contains the user-facing and service-facing binaries:
\begin{itemize}
\item \texttt{source/applications/shell/};
\item \texttt{source/applications/worker/};
\item \texttt{source/applications/gui/};
\item \texttt{source/applications/modelSpecific/}; and
\item the separate \texttt{source/applications/mcp/} workspace.
\end{itemize}
The MCP directory is useful integration material, but it is not added by the
root CMake graph and should therefore be treated as a companion workspace
rather than a compiled target family.

The tools layer is a shared backend library tree in \texttt{source/tools/}.
Its source families currently include AI assistant support, biochemical,
continuous, factorial design, optimization, and statistics. Only
\texttt{source/tools/AIAssistant/} has its own subdirectory CMake file; the
other tool subdirectories are source families consumed through the top-level
\texttt{genesys\_tools} library.

\section{Non-build support areas}
\label{sec:repo-support-areas}

Several top-level directories are part of the repository but are not source
architecture in the narrow CMake sense.

\begin{table}[htbp]
\centering
\caption{Support areas that are part of the repository but not compiled as core architecture.}
\label{tab:repo-support-areas}
\small
\begin{tabular}{|p{3.2cm}|p{4.5cm}|p{4.5cm}|}
\hline
\textbf{Area} & \textbf{Current role} & \textbf{Build graph status} \\
\hline
\texttt{docs/} & Manuals, governance, architectural references, and generated PDFs & documentation source, not runtime architecture \\
\hline
\texttt{models/} & Example and didactic \texttt{.gen} models used by the manual and validation flows & data/content root, not a compiled target family \\
\hline
\texttt{debian/} & Debian packaging metadata and install manifests & packaging-only \\
\hline
\texttt{packaging/} & Linux desktop, Docker, and OVA helpers & packaging-only \\
\hline
\texttt{scripts/} & Maintenance and automation scripts & tooling-only \\
\hline
MCP workspace & Python MCP workspace for assistant integration & integration workspace; not in the root CMake graph \\
\hline
\end{tabular}
\end{table}

The support areas matter because they shape how the code is delivered and
documented, but they do not define the runtime module boundaries of the
simulator itself.

\section{Dependency and build flow}
\label{sec:repo-dependency-build-flow}

The current build flow is a configure-build-test pipeline with separate preset
families for unit tests, shell, worker, GUI, and model-specific applications.
The safest reading is:
\begin{enumerate}
\item configure with the correct preset;
\item build the matching target family;
\item run the matching tests when the preset family defines them;
\item inspect packaging or installation artifacts only after the compiled
surface is confirmed.
\end{enumerate}

\begin{figure}[htbp]
\centering
% FIGURE-SPEC-BEGIN
% ID: fig:repo-boundary-map
% Source of truth:
%   - ../CMakeLists.txt
%   - ../source/applications/gui/CMakeLists.txt
%   - ../source/applications/modelSpecific/CMakeLists.txt
%   - ../source/applications/worker/CMakeLists.txt
%   - ../source/tools/CMakeLists.txt
%   - ../source/tests/CMakeLists.txt
% Purpose:
%   Show which parts of the source tree are in the current compile boundary and which are support-only.
% Required visual content:
%   - source subtrees that compile into targets
%   - support-only roots such as docs/models/packaging/scripts
%   - explicit note that the MCP workspace is outside the root build graph
% Accessibility description:
%   A boundary map that prevents readers from assuming every directory is a current compiled module.
% Validation:
%   The separation between compiled targets and support/workspace roots must match the current root CMake files.
% FIGURE-SPEC-END
\figureplaceholder{figs/generated/repository-boundary-map}
\caption{Current compiled boundary versus support-only repository areas.}
\label{fig:repo-boundary-map}
\end{figure}

Figure~\ref{fig:repo-boundary-map} is the easiest way to avoid false
assumptions. A directory existing on disk does not mean it is in the root build
graph. The compiled boundary is created by \texttt{add\_subdirectory()} calls in
the root CMake file and by the options that enable or disable them.

\begin{figure}[htbp]
\centering
% FIGURE-SPEC-BEGIN
% ID: fig:repo-build-flow
% Source of truth:
%   - ../CMakePresets.json
%   - ../README.md
%   - ../docs/ai_assistants/runbooks/LOCAL_AGENT.md
%   - ../docs/developers/ManualGenESyS_source/make.sh
% Purpose:
%   Show the current configure -> build -> test flow and where packaging/doc generation sits relative to it.
% Required visual content:
%   - preset selection
%   - configure step
%   - build step
%   - test step
%   - packaging/documentation side paths
% Accessibility description:
%   A flow diagram that explains how a developer moves from preset selection to a validated target.
% Validation:
%   The steps and preset names must match the current preset inventory and the documented runbooks.
% FIGURE-SPEC-END
\figureplaceholder{figs/generated/repository-build-flow}
\caption{Current configure, build, test, and support workflow.}
\label{fig:repo-build-flow}
\end{figure}

The current preset inventory makes this flow concrete. For example:
\begin{itemize}
\item \texttt{tests-unit}, \texttt{tests-kernel-unit}, and
\texttt{tests-smoke} are the test-oriented families;
\item \texttt{genesys\_shell} and the model-specific presets build the terminal
family;
\item \texttt{genesys\_worker\_app} and \texttt{gui-app} build the service and
main GUI families;
\item the independent GUI presets build their own executables in separate build
trees.
\end{itemize}

The build tree separation is part of the architecture. It prevents one family
from overwriting another family accidentally and makes it possible to reason
about target boundaries from the preset name alone.

\section{Reading order and final notes}
\label{sec:repo-reading-order}

The most reliable reading order for this repository is:
\begin{enumerate}
\item root \texttt{README.md};
\item \path{docs/ai_assistants/ARCHITECTURE.md};
\item root \texttt{CMakeLists.txt} and \texttt{CMakePresets.json};
\item the \texttt{source/kernel/} subtree;
\item the parser and plugins subtrees;
\item the application and tools subtrees;
\item the tests subtree;
\item the support roots (\texttt{docs/}, \texttt{models/}, \texttt{debian/},
\texttt{packaging/}, and \texttt{scripts/}).
\end{enumerate}

This chapter deliberately avoids overclaiming the current maturity of any
individual application or scientific feature. It is a map of current structure,
not a promise that every directory is equally production-ready.

\section{Test your understanding}
\label{sec:repo-test-your-understanding}

\begin{enumerate}
\item Why is the root \texttt{CMakeLists.txt} a better source of truth for build
architecture than directory names alone?
\item Which directories are support areas rather than compiled runtime
architecture?
\item Why should \texttt{source/applications/mcp/} be treated differently from
the current compiled application families?
\item Which target family forms the runtime core that the applications and tests
share?
\end{enumerate}
